\chapter{Python Spectral Analysis}

This appendix details the Python digital signal processing (DSP) pipeline executed on the host machine. It governs the real-time serial buffering and the mathematical transformation of the physical acoustic data.

\section{Real-Time Buffering and STFT Spectrogram}

The following script utilizes the \texttt{pyserial}, \texttt{numpy}, and \texttt{scipy} libraries to interface with the ESP32. It buffers the incoming 12-bit audio data, applies a mathematical Hanning window, and computes the Short-Time Fourier Transform (STFT) to generate a dynamic spectrogram of the cardiopulmonary acoustic sweep.

\begin{lstlisting}[language=Python, caption={Python DSP script for real-time serial data buffering and STFT spectrogram generation.}, label={lst:python_dsp}]
import serial
import numpy as np
import matplotlib.pyplot as plt
from scipy.signal import spectrogram

# Configure Serial Port (Update 'COM3' or '/dev/ttyUSB0' as needed)
SERIAL_PORT = 'COM3'
BAUD_RATE = 500000
FS = 20000  # 20 kHz sampling rate

def acquire_data(duration_seconds=5):
    """Buffers incoming UART data from the ESP32."""
    samples_to_read = FS * duration_seconds
    audio_data = np.zeros(samples_to_read)
    
    with serial.Serial(SERIAL_PORT, BAUD_RATE, timeout=1) as ser:
        print(f"Recording {duration_seconds} seconds of acoustic data...")
        for i in range(samples_to_read):
            try:
                line = ser.readline().decode('utf-8').strip()
                audio_data[i] = int(line)
            except ValueError:
                audio_data[i] = audio_data[i-1] # Hold previous value on error
                
    # Center waveform around 0
    audio_data = audio_data - np.mean(audio_data)
    return audio_data

def plot_spectrogram(data):
    """Computes and visualizes the STFT Spectrogram."""
    f, t, Sxx = spectrogram(data, fs=FS, window='hann', nperseg=1024, noverlap=512)
    
    plt.figure(figsize=(10, 6))
    # Plot spectrogram in dB
    plt.pcolormesh(t, f, 10 * np.log10(Sxx + 1e-10), shading='gouraud', cmap='jet')
    plt.ylim(0, 2000) # Limit view to cardiopulmonary bandwidth
    plt.ylabel('Frequency [Hz]')
    plt.xlabel('Time [sec]')
    plt.title('Cardiopulmonary Acoustic Spectrogram (20 Hz - 2000 Hz)')
    plt.colorbar(label='Intensity [dB]')
    plt.show()

if __name__ == '__main__':
    raw_audio = acquire_data(duration_seconds=10)
    plot_spectrogram(raw_audio)
\end{lstlisting}